\section{结论}

本文面向区域风电场群 24 步集群总功率预测问题，研究了如何在严格时间可用性约束下联合利用近期场站级功率和未来网格化气象预报。本文的核心思路是对两类输入赋予不同角色：未来气象预报提供预测窗口内的主要外生驱动，近期功率记录表征目标窗口前的当前运行状态、场站间耦合和短期惯性。基于这一设定，本文提出功率条件化气象记忆网络，将近期功率动态编码为样本特定的 dynamic power context tokens，并注入气象记忆空间，使 future queries 能够在状态条件化后的气象记忆中读取预测信息。

吉林区域风电场群案例研究表明，所提出模型在整体 24 步预测任务上优于 Weather-only 以及多种网格气象编码基线，并在 MAE、RMSE 和 \(R^2\) 上取得更优结果。逐步长分析显示，近期功率上下文的贡献主要集中在预测窗口前段，这与其刻画当前运行状态和短期惯性的作用一致；远端 horizon 上模型表现则更多受未来气象预报质量影响。

消融实验进一步验证了模型设计的有效性。真实功率输入、多尺度图功率编码、气象空间记忆、气象变化分支和变量解耦编码均对预测性能产生贡献；气象通道消融也表明，风场相关变量是该任务中不可替代的关键输入，背景气象变量主要提供补充性约束。

总体而言，本文结果说明，在区域风电集群总功率预测中，将近期功率作为动态状态上下文、将未来气象作为主要驱动信息，是一种有效且符合业务时间约束的建模方式。未来工作可进一步围绕跨区域迁移、气象预报不确定性建模，以及极端爬坡、限电恢复和概率预测等调度相关场景展开。
